% sec:fd-orders — Second- and Fourth-Order Differences
\section{Second- and Fourth-Order Differences}
\label{sec:fd-orders}

The Vandermonde system of \cref{sec:fd-stencils} produces stencils of any
order, but two choices dominate numerical relativity: second order ($m = 1$,
three-point stencils) and fourth order ($m = 2$, five-point stencils).  The
decision between them is a trade-off: fourth-order stencils cost more
memory accesses per evaluation but converge so much faster that they require
far fewer grid points to reach a given accuracy.  Quantifying that trade-off
requires the explicit stencil formulae, their leading error coefficients,
and a direct numerical comparison.

\paragraph{The four centred stencils.}

The second-order first derivative was derived in
\cref{eq:d1-second-order}; the fourth-order first and second derivatives
in \cref{eq:d1-fourth-order,eq:d2-fourth-order}.  The remaining member of
the set is the second-order centred second derivative, obtained from the
Vandermonde system with $m = 1$, $p = 2$:
%
\begin{equation}
  f''_i \approx \frac{f_{i+1} - 2\,f_i + f_{i-1}}{\Delta x^2} \,.
  \label{eq:d2-second-order}
\end{equation}
%
The three-point stencil is the simplest centred approximation to $f''$: it
measures how much the central value $f_i$ deviates from the average of its
two neighbours --- exactly the discrete analogue of curvature.

\paragraph{Leading error coefficients.}

Each stencil's truncation error is determined by the first uncancelled term
in the Taylor expansion (\cref{eq:taylor-grid}).  Evaluating the coefficient
of that term gives the leading error explicitly.  For the four centred
stencils:

\vspace{0.5\baselineskip}
\noindent\makebox[\textwidth][c]{\begin{tabular}{llccl}
  \toprule
  Derivative & Order & Points & Stencil & Leading error \\
  \midrule
  $f'_i$ & 2nd & 3 &
    $\dfrac{f_{i+1} - f_{i-1}}{2\,\Delta x}$ &
    $\dfrac{\Delta x^2}{6}\,f'''_i$ \\[10pt]
  $f'_i$ & 4th & 5 &
    $\dfrac{-f_{i+2} + 8f_{i+1} - 8f_{i-1} + f_{i-2}}{12\,\Delta x}$ &
    $\dfrac{\Delta x^4}{30}\,f^{(5)}_i$ \\[10pt]
  $f''_i$ & 2nd & 3 &
    $\dfrac{f_{i+1} - 2f_i + f_{i-1}}{\Delta x^2}$ &
    $\dfrac{\Delta x^2}{12}\,f^{(4)}_i$ \\[10pt]
  $f''_i$ & 4th & 5 &
    $\dfrac{-f_{i+2} + 16f_{i+1} - 30f_i + 16f_{i-1} - f_{i-2}}
    {12\,\Delta x^2}$ &
    $\dfrac{\Delta x^4}{90}\,f^{(6)}_i$ \\
  \bottomrule
\end{tabular}\hspace*{30mm}}

\vspace{0.5\baselineskip}
\noindent The error coefficients follow from the Vandermonde analysis: the $n$-th
uncancelled term contributes $\sum_j w_{k_j}\,k_j^n / n!$ times
$\Delta x^{n-p}\,f^{(n)}_i$.  For centred stencils the first uncancelled
term after the $N = 2m + 1$ conditions is the one whose parity matches
the derivative order, so the leading exponent is always $\Delta x^{2m}$.
\physnote{The error coefficients decrease as the stencil order increases:
$1/6 \to 1/30$ for the first derivative, $1/12 \to 1/90$ for the second.
This means that even at the same grid spacing, the fourth-order stencil has
a smaller prefactor in addition to its faster convergence rate.}

\paragraph{A worked comparison.}

To see what the error coefficients mean in practice, consider the test
function $f(x) = \sin(2\pi x)$ on the interval $[0, 1]$ with $N_x = 20$
grid points ($\Delta x = 0.05$).  The first derivative at $x = 0.5$ is
$f'(0.5) = 2\pi\cos(\pi) = -2\pi \approx -6.2832$.  The exact higher
derivatives needed for the error terms are
$f'''(0.5) = -(2\pi)^3\cos(\pi) = 8\pi^3$
and $f^{(5)}(0.5) = (2\pi)^5\cos(\pi) = -32\pi^5$.

The second-order stencil gives an error of approximately
%
\begin{equation}
  \frac{\Delta x^2}{6}\,\abs{f'''(0.5)}
  = \frac{(0.05)^2}{6} \cdot 8\pi^3
  \approx 1.03 \times 10^{-1} \,,
  \label{eq:d1-error-second-order}
\end{equation}
%
while the fourth-order stencil gives
%
\begin{equation}
  \frac{\Delta x^4}{30}\,\abs{f^{(5)}(0.5)}
  = \frac{(0.05)^4}{30} \cdot 32\pi^5
  \approx 2.04 \times 10^{-3} \,.
  \label{eq:d1-error-fourth-order}
\end{equation}
%
The fourth-order stencil is roughly $50$ times more accurate, at the cost
of accessing two additional grid points per evaluation ($5$ versus $3$).
Doubling the resolution to $N_x = 40$ reduces the second-order error by a
factor of $4$ (to $\sim 2.6 \times 10^{-2}$) and the fourth-order error by
a factor of $16$ (to $\sim 1.3 \times 10^{-4}$).  The gap widens with each
refinement: to match the fourth-order accuracy at $N_x = 20$, the
second-order scheme needs roughly $N_x \approx 144$ --- over seven times as
many points per direction, and about $370$ times the three-dimensional grid
volume.
\implnote{The codebase's convergence tests verify these scaling predictions
numerically.  The test \texttt{test\_d1x\_sin\_convergence} in
\codemodule{stencils} applies the fourth-order stencil to $\sin(2\pi x)$
at resolutions $N = 16, 32, 64$ and checks that the error ratio between
successive doublings lies in $[12, 20]$ --- bracketing the theoretical
value of $16$.}

\paragraph{Cost-benefit analysis.}

The computational cost of a stencil evaluation scales with the number of
grid-point accesses, not with arithmetic operations --- modern processors
can multiply and add far faster than they can fetch data from main memory.
A three-point stencil touches $3$ values; a five-point stencil touches $5$.
In three spatial dimensions, the BSSN right-hand side requires first and
second derivatives in each of three directions plus mixed second derivatives
for all three coordinate pairs, for a total of $3 + 3 + 3 \times 2 = 12$
stencil evaluations per field per grid point (counting each pass of the
separable mixed-derivative computation separately).

At second order, each evaluation touches $3$ points, but many accesses
overlap --- the same grid value appears in multiple stencils.  At fourth
order the overlap is even greater because the five-point stencil in one
direction shares its central three points with the three-point stencil of
the lower order.  In practice, the actual cost increase from second to
fourth order is modest: the codebase measures roughly a $30$--$50\%$
increase in wall-clock time per right-hand-side evaluation, depending on
cache behaviour.

The accuracy gain dwarfs this cost.  Reaching a target error of $10^{-6}$
with a second-order scheme requires roughly $\Delta x \sim 10^{-3}$, or
about $10^3$ points per direction --- a billion grid points in three
dimensions.  The fourth-order scheme reaches the same error with
$\Delta x \sim 10^{-3/2} \approx 0.03$, or about $30$ points per direction
--- roughly $27\,000$ grid points.  The memory ratio is $10^9 / 2.7 \times
10^4 \approx 37\,000$, and the compute-time ratio is even larger because
the time integrator's stability constraint (the CFL condition, discussed in
\cref{ch:ode-integration}) forces smaller time steps at finer resolution.
\warnnote{These estimates assume the solution is smooth enough that the
truncation error dominates.  Near discontinuities or sharp features, the
higher-order stencil can produce oscillations (Gibbs phenomenon) that the
lower-order stencil avoids.  Kreiss--Oliger dissipation
(\cref{sec:fd-dissipation}) addresses this by selectively damping the
grid-scale modes that trigger oscillations.}

For the smooth fields of numerical relativity --- conformal factor
$\conformal$, extrinsic curvature $\extrinsic$, lapse $\lapse$, shift
$\shift$ --- the fourth-order stencil is the clear
winner, and the codebase uses it exclusively for interior spatial
derivatives.  The second-order stencil appears only at boundaries and in
the pedagogical comparison above.
